\section{CONCLUSIONS}

The present study successfully analyzed the dynamic behavior of periodic structures applied to TLP tendons, demonstrating the efficacy of geometric and material periodicity as passive vibration control strategies. By modeling tendons as EBT beams, configured as both stepped and bi-material structures, a clear relationship between the unit cell geometry and the formation of frequency band gaps was established.

The implementation of FEM routines in MATLAB yielded consistent results when validated against analytical models, achieving a maximum percentile error of approximately 0.000688\% for the first ten natural frequencies.

The introduction of the geometric saturation parameter $\alpha$, proved to be a critical design variable for tuning attenuation zones. The parametric analysis revealed distinct behaviors for the two configurations investigated:

\begin{itemize}
    \item Bi-material configuration: While capable of generating band gaps, such as the gap observed between modes $20$ and $21$, the bi-material beam showed limited tunability. The attenuation zones exhibited significant spectral migration, often traversing across mode indices as the saturation parameter varied, or required specific saturation parameters (\textit{e.g.}, near $\alpha=8/16$) to manifest effectively.
    \item Stepped configuration: Introducing geometric periodicity resulted in superior vibration mitigation capabilities. This beam type exhibited multiple, wider band gaps, notably a maximum gap observed at $\alpha=11/16$. Analysis confirmed that geometric variations enable the creation of two distinct attenuation zones across diverse configurations, rendering the stepped beam a more tunable and efficient structure for TLP applications than its bi-material counterpart.
\end{itemize}

In summary, optimizing the saturation parameter $\alpha$ for a stepped beam configuration offers a promising engineering solution for passive mitigation of vibration in offshore platforms. These findings contribute to the development of safer and more durable marine structures by effectively filtering vibrations on revelant frequency spectrum.